\subsection{两个函数的复合。反函数}

命题6. 若 \(f\) 是从A到B的映射，且 \(g\) 是从B到C的映射，则 \(g \circ f\) 是从A到C的映射。

设F、G分别为f、g的图像，我们来证明G∘F是一个函数图像。设x、z、 \(z'\) 为对象，使得 \((x, z) \in \mathbf{G} \circ \mathbf{F}\) , \((x, z') \in \mathbf{G} \circ \mathbf{F}\) 。存在对象y、 \(y'\) 使得

\[
(x, y) \in \mathrm{F}, \quad (x, y ^ {\prime}) \in \mathrm{F}, \quad (y, z) \in \mathrm{G}, \quad (y ^ {\prime}, z ^ {\prime}) \in \mathrm{G}.
\]

由于 F 是函数图，我们有 \(y = y'\) 因此， \((y, z') \in G\) 。由于 G 是函数图，由此可得 \(z = z'\) ，这证明了我们的断言。此外， \(g \circ f\) 的定义域显然是 A，证明完成。

函数 \(g \circ f\) 也可以写成 \(x \to g(f(x))\) ，或写成 \(gf\) 。

定义 10. 设 f 是从 A 到 B 的映射。若 A 中任意两个不同元素在 f 下的像不同，则称映射 f 是单射的，或称为单射。～若 \(f(A) = B\) ，则称映射 f 是满射的，或称为满射。如果 f 既是单射又是满射，则称其为双射，或一一对应。

我们有时不说 f 是满射，而说 f 是 A 到 B 上的映射，或者说它是 B 借助 A 的参数表示（此时 A 称为该表示的参数集，A 的元素称为参数）。如果 f 是双射，我们有时说 f 使 A 与 B 一一对应。A 到 A 的双射称为 A 的置换。

\minorheading{示例}

\begin{enumerate}
\def\labelenumi{(\arabic{enumi})}
\item
  如果 \(A \subset B\) ，则从 \(A\) 到 \(B\) 且以 \(A\) 的对角线为图的映射是单射，称为 \(A\) 到 \(B\) 的典范映射或典范单射（或简称单射）。
\item
  设 A 是一个集合。映射 \(x \to (x, x)\) 从 A 到 \(\Delta_{A}\) （ \(A \times A\) 的对角线）是双射，称为 A 的对角映射。
\item
  设 G 是一个有序对集合。映射 \(pr_{1}\) （相应地 \(pr_{2}\) ）从 G 到 \(pr_{1}G\) （相应地 \(pr_{2}G\) ）是满射； \(pr_{1}\) 是单射当且仅当 G 是函数图。
\item
  设 G 是一个有序对集合。映射
\end{enumerate}

\[
z \rightarrow (\mathrm{pr} _ {2} z, \mathrm{pr} _ {1} z)
\]

是从 \(\mathbf{G}\) 到 \(\overline{\mathbf{G}}\) 的双射（称为典范双射）。

\begin{enumerate}
\def\labelenumi{(\arabic{enumi})}
\setcounter{enumi}{4}
\tightlist
\item
  设 \(\mathbf{A}\) 是一个集合， \(b\) 是一个对象。映射 \(x \to (x, b)\) 是从 \(\mathbf{A}\) 到 \(\mathbf{A} \times \{b\}\) 的双射。
\end{enumerate}

命题7. 设 \(f\) 为从 A 到 B 的映射。 \(\overline{f}^1\) 是一个函数，当且仅当 \(f\) 是双射的。

如果 \(\overline{f}^{1}\) 是一个函数，其源 B 等于其定义域，即等于 \(f(A)\) ；因此 f 是满射的。为了证明 f 是单射的，设 x 和 y 是 A 的两个元素，使得 \(f(x)=f(y)\) 。若 F 表示 f 的图像，我们有

\[
(f (x), x) \in \overline {{\mathrm{F}}} ^ {- 1} \quad \text { 且 } \quad (f (y), y) \in \overline {{\mathrm{F}}},
\]

因此

\[
(f (x), y) \in \overline {{\mathbf {F}}},
\]

使得 \(x = y\) ，这证明了该断言。

反之，如果 \(f\) 是双射的，则立即可得出 \(\overline{\mathbf{F}}\) 是函数性的，且 \(\overline{f}^1\) 的定义域等于 \(\mathbf{B}\) .

当 \(f\) 是双射时， \(\overline{f}^1\) 称为 \(f\) 的逆映射; \(\overline{f}^1\) 是双射， \(\overline{f}^1 \circ f\) 是 A 的恒等映射，并且 \(f \circ \overline{f}^1\) 是 B 的恒等映射。

如果一个置换与其逆置换相同，则称其为对合的。

备注。设 \(f\) 是从 \(A\) 到 \(B\) 的映射。对于 \(A\) 的每个子集 \(X\) ，都有（第3号） \(X \subset f\langle f(X)\rangle\) 。此外，对于 \(B\) 的每个子集 \(Y\) ，都有 \(f\langle f(Y)\rangle \subset Y\) ，因为关系 \(y \in f\langle f(Y)\rangle\) 等价于

\[
(\exists x) ((\exists z) (z \in Y \text { 且 } z = f (x)) \text { 且 } y = f (x))
\]

因此蕴含关系 \((\exists z)(z \in \mathbf{Y}\) 并且 \(y = z)\) 并且因此也有关系 \(y \in \mathbf{Y}\) .

如果 \(f\) 由于是满射，我们有 \(f\langle f\langle \mathrm{Y}\rangle\rangle = \mathrm{Y}\) 对于每一个子集 \(\mathrm{Y}\) 的 \(\mathrm{B}\) ，对于关系 \(y\in \mathrm{Y}\subset \mathrm{B}\) 由假设推出关系 \((\exists x)(y = f(x))\) ，因此也 \((\exists x)(y\in \mathrm{Y}\) 并且 \(y = f(x)\) ）；但“ \(y\in \mathrm{Y}\) 并且 \(y = f(x)\) '' 蕴含 \((\exists z)(z\in \mathrm{Y}\) 并且 \(z = f(x)\) ），于是我们的断言成立。

如果 \(f\) 是单射，我们有 \(f\langle f\langle \mathbf{X}\rangle = \mathbf{X}\) 对于每一个子集 \(\mathbf{X}\) 的 \(\mathbf{A}\) 对于关系 \(x\in f^{-1}\langle f\langle \mathbf{X}\rangle\rangle\) 等价于 \(f(x)\in f\langle \mathbf{X}\rangle\) ，因此到

\[
(\exists z) (z \in \mathrm{X} \text {   且   } f (z) = f (x));
\]

但假设意味着 \(f(z) = f(x)\) 蕴含 \(z = x\) ，所以 \(x \in f\langle f\langle \mathbf{X}\rangle \rangle\) 蕴含 \(x \in \mathbf{X}\) .

